% GENERATED_BY: run_oc_core_monograph_translation_rework_dispatch_v1.py
% AI_ASSISTED_TRANSLATION_DRAFT: TRUE
% THEOREM_REVIEW_STATUS: APPROVED
% THEOREM_REVIEW_MODE: FORMAL_AI_GATE__CODEX_CLI_CHATGPT
% THEOREM_REVIEW_REVIEWER_ID: CODEX_CLI_CHATGPT__AUTO_DISPATCH_20260406
% THEOREM_REVIEW_AUTHORITY_CLASS: FINAL_ADVERSARIAL_EXTERNAL_LLM_REVIEWER
% THEOREM_REVIEWED_AT: 2026-04-14T11:09:57Z
% SOURCE_LANGUAGE: EN
% TARGET_LANGUAGE: RU
% SOURCE_EN_REF: appendix/E_oc_core_1_3_journal_core_bridge.tex
\section{Мост извлечения журнального ядра}

Это приложение объясняет, как статья журнального ядра соотносится с основной монографией.

\subsection{Зачем существует журнальное ядро}

Некоторым каналам требуется ограниченный по объему, более быстрый для чтения научный артефакт.
Редакторам, рецензентам и внешним читателям часто нужен объект формата статьи, чье защищаемое утверждение можно
прочитать за один присест.
Журнальное ядро существует именно для этой цели.
Его задача не в том, чтобы заменить монографию, а в том, чтобы предоставить узкую статью, которая остается привязанной к
тем же научным условиям истинности.

\subsection{Что журнальное ядро может опускать}

Журнальное ядро может сжимать:

\begin{itemize}
\item развернутое дидактическое объяснение на уровне глав,
\item большие части атласа предметной области,
\item широкий охват модулей,
\item подробные инвентаризации приложений,
\item а также часть педагогического каркаса уровня монографии.
\end{itemize}

Оно не может без явного указания изменять:

\begin{itemize}
\item объем теоремы,
\item явно указанные положения, которые не утверждаются,
\item заявленную базу источников,
\item смысл рисунков,
\item или связь между унаследованным корпусом и защищаемым здесь утверждением.
\end{itemize}

\subsection{Что монография гарантирует последующим артефактам}

Поскольку журнальное ядро извлечено из настоящего тома, последующим каналам, адресованным владельцу, и
публичным каналам не нужно импровизировать собственное понимание науки.
Они могут извлекать материал из основной монографии, когда им нужна глубина, из журнального ядра, когда им нужно
ограниченное по объему изложение, и из сопроводительного слоя, когда им нужны подтверждающие материалы.

В этом состоит практическая публикационная ценность монографии.
Она не просто объемнее статьи.
Это документ, который не дает каждой последующей проекции превратиться в отдельное интерпретационное
упражнение.

\clearpage
